% =============================================================================
% Section VIII: Conclusions
% paper/conclusions.tex
% =============================================================================

\section{Conclusions}
\label{sec:conclusions}

We have computed the decoherence number $\Nnum$ for a charged body in spatial
superposition in the family of spatially flat FLRW spacetimes with power-law
scale factor $\Scfac(\tau) = C_{p}(\tau - \tau_{*})^{p}$, $p > 1$, using the
DSW local decoherence framework with the conformal-vacuum Wightman function.
The key input is that conformal invariance of the massless scalar (and Maxwell)
field causes the scale-factor dependence to cancel exactly from the decoherence
integrand, reducing the problem to a purely flat-space double integral over
conformal time.  The closed-form result,
%
\begin{equation}
  \Nnum(T)
  =
  \frac{q^{2}d^{2}p^{3}}{48\pi^{3}}
  \cdot
  \frac{2T_{0}^{2} + \widetilde{T}(2 + \widetilde{T})}{T_{0}^{2}(1 + \widetilde{T})^{2}}
  \cdot
  \ln(1 + \widetilde{T}) \,,
  \label{eq:N-conclusions}
\end{equation}
%
where $\widetilde{T} = T/T_{0}$ and $T_{0}$ is the proper-time scale set by
the initial conditions, holds for any $T > 0$ and any $p > 1$.  For
$T \gg T_{0}$ this grows as $\Nnum \sim \ln T$, diverging as $T \to \infty$.
Decoherence is \emph{inevitable}: no matter how well isolated the superposition,
it is fully decohered at sufficiently late times.  This conclusion holds for the
entire power-law family, none of which possesses a Killing symmetry or a
thermal structure analogous to the Gibbons-Hawking temperature.

The causal horizon is identified as the essential mechanism by the sharp
contrast with the $p = 1$ case, in which no future causal horizon is
present and $\Nnum$ saturates to a finite constant as $T \to \infty$: the
off-diagonal coherence does not vanish.  The transition from saturation to
logarithmic divergence at $p = 1$ demonstrates that it is the irreversible
sequestration of soft radiation beyond the horizon — not the dynamics of
the scale factor per se — that converts bounded decoherence into unbounded
decoherence.  Specializing to de Sitter spacetime, \cref{eq:N-conclusions}
reduces exactly to the known result
$\Nnum_{\rm dS} = q^{2}d^{2}H^{3}T/(96\pi^{2})$~\cite{Danielson:2022tdw},
recovering both the linear growth in $T$ and the numerical coefficient.
The linear rather than logarithmic growth in de Sitter reflects the
exponential expansion's constant rate of sweeping modes through the Killing
horizon; written in terms of the Gibbons-Hawking temperature $T_{\rm GH} =
H/2\pi$, the rate $\Nnum_{\rm dS}/T \propto T_{\rm GH}^{3}$ has the
characteristic form of thermal decoherence by a heat bath, a structure
that is absent in the non-stationary power-law family.

Conformal invariance gives the electromagnetic result without further
computation: the two independent transverse polarizations of the photon
contribute equally, yielding $\Nnum_{\rm EM} = 2\,\Nnum_{\rm scalar}$
for every conformally flat FLRW background, with no dependence on the
scale factor or on whether the horizon is a Killing horizon.  This
universality is a direct kinematic consequence of the field's conformal
structure: the polarization count is background-independent.  Taken
together, the scalar and electromagnetic results establish a general
principle: any conformally invariant massless field in an accelerating
FLRW cosmology with a future causal horizon will inevitably decohere
spatial superpositions of a coupled source, with a rate that grows at
least logarithmically in the superposition time.

The most pressing open direction is the extension to soft gravitons in
FLRW spacetimes.  The linearized gravitational field is not conformally
invariant, so the methods developed here do not apply directly; new
techniques are required to handle the non-conformal mode structure and
the gauge redundancy of the graviton.  The present results provide a
clean semiclassical baseline — exact decoherence rates from conformally
invariant fields — against which future quantum-gravity investigations of
horizon decoherence can be calibrated.

%% CITATIONS NEEDED (for gpd-bibliographer):
%% - Danielson:2022tdw  -- DSW (2023) PRD 108, 025007, arXiv:2301.00026
